%%%%%%%%%%%%%%%%%%%%%%%%%%%%%%%%%%%%%%%%%%%%%%%%%%%%%%%%%%%%
%% AUFGABE 4
%%%%%%%%%%%%%%%%%%%%%%%%%%%%%%%%%%%%%%%%%%%%%%%%%%%%%%%%%%%%

\newpage
\section{Wertebestimmung von Widerstand und Kondensator}

\subsection{Vorbereitung}
\subsubsection{Widerstandsbestimmung}

\newcommand{\resistorRow}[3]{%
    \noindent
    \begin{minipage}[c]{0.4\textwidth}
        \includegraphics[width=5cm]{#1} 
    \end{minipage}%
    \hfill
    \begin{minipage}[c]{0.5\textwidth}
        \large
        Widerstandswert: \underline{\hspace{0.2cm} #2 \hspace{0.2cm}} \\[0.3cm]
        Toleranz: \hspace{1.3cm} \underline{\hspace{0.2cm} #3 \hspace{0.2cm}}
    \end{minipage}
    
    \vspace{1cm} % Abstand zur nächsten Reihe
}

\resistorRow{resistor1.png}{\SI{47}{\kilo\ohm}}{$\pm \SI{5}{\percent}$}

\resistorRow{resistor2.png}{\SI{3,9}{\mega\ohm}}{$\pm \SI{5}{\percent}$}

\resistorRow{resistor3.png}{\SI{6,8}{\kilo\ohm}}{$\pm \SI{10}{\percent}$}

\subsubsection{Spannungsteiler}

\paragraph{a)}
Die Spannung teilt sich am Spannungsteiler im Verhältnis der Widerstände auf.
\[
    U_{R1} = U_e - U_a = \SI{5}{\volt} - \SI{3}{\volt} = \SI{2}{\volt}
\]
Das Verhältnis der Widerstände $R_1 : R_2$ entspricht dem Verhältnis der Spannungen $U_{R1} : U_a$, also \textbf{2:3}.

\noindent Berechnung von $R_2$ mit $R_1 = \SI{10}{\kilo\ohm}$:
\begin{align*}
    \frac{R_1}{R_2} &= \frac{U_{R1}}{U_a} \\
    R_2 &= R_1 \cdot \frac{U_a}{U_{R1}} \\
    R_2 &= \SI{10}{\kilo\ohm} \cdot \frac{\SI{3}{\volt}}{\SI{2}{\volt}} \\
    R_2 &= \SI{15}{\kilo\ohm}
\end{align*}

\paragraph{b)}
Der Widerstand $R_2$ wird durch $R_x$ ersetzt. Die neue Ausgangsspannung beträgt $U_a = \SI{1}{\volt}$.
\[
    U_{R1} = \SI{5}{\volt} - \SI{1}{\volt} = \SI{4}{\volt}
\]
Berechnung von $R_x$:
\begin{align*}
    R_x &= R_1 \cdot \frac{U_a}{U_{R1}} \\
    R_x &= \SI{10}{\kilo\ohm} \cdot \frac{\SI{1}{\volt}}{\SI{4}{\volt}} \\
    R_x &= \SI{2.5}{\kilo\ohm}
\end{align*}

\subsubsection{Kondensatorladung}

\paragraph{a)} Die Spannung am Kondensator folgt der Ladefunktion:
\[
u_c(t) = U_q \cdot \left(1 - e^{-\frac{t}{\tau}}\right)
\]
Es ist bekannt, dass der Kondensator nach der Zeitkonstante $\tau = R \cdot C$ circa \SI{63}{\percent} der Endspannung erreicht hat ($1 - e^{-1} \approx 0.63$).

\noindent Da nach $t = \SI{100}{\milli\second}$ genau dieser Wert erreicht ist, gilt:
\[ \tau = \SI{100}{\milli\second} \]

\noindent Daraus lässt sich die Kapazität $C$ berechnen:
\begin{align*}
    \tau &= R \cdot C \\
    C &= \frac{\tau}{R} \\
    C &= \frac{\SI{0.1}{\second}}{\SI{10}{\kilo\ohm}} \\
    C &= \SI{0.00001}{\farad} = \SI{10}{\micro\farad}
\end{align*}

\paragraph{b)} Der Kondensator gilt als vollständig geladen (Spannung $> 99\,\%$), wenn eine Zeit von ca. $5 \cdot \tau$ vergangen ist ($1 - e^{-5} \approx 0.993$).

\noindent Die Ladedauer beträgt somit:
\begin{align*}
    t_{\text{voll}} &\approx 5 \cdot \tau \\
    t_{\text{voll}} &= 5 \cdot \SI{100}{\milli\second} \\
    t_{\text{voll}} &= \SI{500}{\milli\second}
\end{align*}

\subsection{Durchführung}

\begin{itemize}
    \item \href{https://www.tinkercad.com/things/6dP6bOFcvVJ-4-widerstandsmessung?sharecode=cloxrV9fQVUj9Cfl4C-4CEYhAw4g13SRQi1Q5vM9bJg}{Tinker CAD Link Aufgabe 1}
    \item \href{https://www.tinkercad.com/things/84XtVuCZHAm-4-kapazitatsmessung?sharecode=4kCbAaRwm3v1nCWt6jAq8217Vbh7pe0KWzrzvLyUI0Q}{Tinker CAD Link Aufgabe 2}
    \item Stückliste
\end{itemize}

\newpage
\subsection{Ergebnisse}

\subsubsection{Widerstand 1}

\paragraph{Farbcode}

Braun Schwarz Schwarz Schwarz Braun

\SI{100}{\ohm}

\paragraph{Mulitmeter}

gemessen: \SI{100.3}{\ohm} $\pm$ \SI{1}{\percent}

\paragraph{UI-Messung}

\begin{table}[h]
    \caption{}\label{}
\begin{tabularx}{\textwidth}{YYY}
    \toprule
Strom / \si{\ampere} & Spannung / \si{\volt} & Widerstand / \si{\ohm} \\ \midrule
0.500 & 0.0505 & \\ \bottomrule
\end{tabularx}
\end{table}

\paragraph{Programm}

hHalo

\begin{table}[h]
    \caption{}\label{}
\begin{tabularx}{\textwidth}{YY}
\toprule
Widerstandswert / \si{\ohm} & ADC Messwert \\ \midrule
78.50 & 8 \\ 
58.76 & 6 \\
68.62 & 7 \\
78.50 & 8 \\ \bottomrule
\end{tabularx}
\end{table}

\subsubsection{Widerstand 2}

\paragraph{Mulitmeter} gemessen: \SI{12.01}{\kilo\ohm}

\paragraph{Farbcode} 

4 Ringe

Braun Rot Orange Gold

\SI{12}{\kilo\ohm} $\pm$ \SI{5}{\percent}

\begin{table}[h]
    \caption{}\label{}
\begin{tabularx}{\textwidth}{YYY}
    \toprule
Strom / \si{\ampere} & Spannung / \si{\volt} & Widerstand / \si{\ohm} \\ \midrule
0.231 & 2.759 & \\ \bottomrule
\end{tabularx}
\end{table}

\paragraph{Programm}

\begin{table}[h]
    \caption{}\label{}
\begin{tabularx}{\textwidth}{YY}
\toprule
Widerstandswert / \si{\ohm} & ADC Messwert \\ \midrule
11891.41 & 557 \\
11891.41 & 557 \\
11891.41 & 557 \\
11891.41 & 557 \\ \bottomrule
\end{tabularx}
\end{table}

\subsubsection{Widerstand 3}

\paragraph{Multimeter} gemessen: \SI{185.7}{\kilo\ohm}

\paragraph{Farbcode} Der Widerstand weist 5 Ringe mit folgenden Farben auf:

\begin{tabularx}{\textwidth}{Y|Y|Y|Y|Y}
    \color{white}\textbf{Braun}\cellcolor{brown} &
    \color{white}\textbf{Grau}\cellcolor{gray} &
    \color{white}\textbf{Schwarz}\cellcolor{black} &
    \color{black}\textbf{Orange}\cellcolor{orange} &
    \color{white}\textbf{Braun}\cellcolor{brown}
\end{tabularx}

Diese Farbkombination entspricht \SI{180}{\kilo\ohm} $\pm$ \SI{1}{\percent}.

\paragraph{Programm} Vom Programm wurden folgende Werte über die Serielle Schnittstelle empfangen:

\begin{table}[h]
    \caption{}\label{}
\begin{tabularx}{\textwidth}{YY}
    \toprule
Widerstandswert / \si{\ohm} & ADC Messwert \\ \midrule
179090.73 & 970 \\
186362.31 & 972 \\
182657.92 & 971 \\
194215.59 & 974 \\
\bottomrule
\end{tabularx}
\end{table}

\paragraph{UI-Messung}

\begin{table}[h]
    \caption{}\label{}
\begin{tabularx}{\textwidth}{YYY}
    \toprule
Strom / \si{\milli\ampere} & Spannung / \si{\volt} & Widerstand / \si{\ohm} \\ \midrule
0.026 & 4.8 & \\ \bottomrule
\end{tabularx}
\end{table}


\subsubsection{Widerstand 4}

\paragraph{Ohmmeter}

gemessen: \SI{1.01}{\kilo\ohm}

\paragraph{UI-Messung}

\begin{table}[h]
    \caption{}\label{}
\begin{tabularx}{\textwidth}{YYY}
    \toprule
Strom / \si{\ampere} & Spannung / \si{\volt} & Widerstand / \si{\ohm} \\ \midrule
0.459 & 0.464 & \\ \bottomrule
\end{tabularx}
\end{table}

\paragraph{Farbcode} Der Widerstand weist 5 Ringe mit folgenden Farben auf:

\begin{tabularx}{\textwidth}{Y|Y|Y|Y|Y}
    \color{white}\textbf{Braun}\cellcolor{brown} &
    \color{white}\textbf{Schwarz}\cellcolor{black} &
    \color{white}\textbf{Schwarz}\cellcolor{black} &
    \color{white}\textbf{Braun}\cellcolor{brown} &
    \color{white}\textbf{Braun}\cellcolor{brown}
\end{tabularx}

Diese Farbkombination entspricht \SI{1}{\kilo\ohm} $\pm$ \SI{1}{\percent}.

\paragraph{Programm}

HAllo

\begin{table}[h]
    \caption{}\label{}
\begin{tabularx}{\textwidth}{YY}
    \toprule
Widerstandswert / \si{\ohm} & ADC Messwert \\ \midrule
960.71 & 90 \\
972.42 & 91 \\
960.71 & 90 \\
960.71 & 90 \\ \bottomrule
\end{tabularx}
\end{table}

\subsubsection{Kondensator 1}

\paragraph{Marking code} Am Bauteil wurde der Wert \SI{220}{\micro\farad} abgelesen.

\paragraph{Programm} Vom Programm wurden folgende Werte über die Serielle Schnittstelle empfangen:

\begin{itemize}
    \item Charge Time: \SI{2397484}{\micro\second}
    \item Capacity: \SI{239.75}{\micro\farad}
\end{itemize}

\subsubsection{Kondensator 2}

\paragraph{Marking code} Am Bauteil wurde der Wert \SI{47}{\micro\farad} abgelesen.

\paragraph{Programm} Vom Programm wurden folgende Werte über die Serielle Schnittstelle empfangen:

\begin{itemize}
    \item Charge Time: \SI{449468}{\micro\second}
    \item Kapazität: \SI{44.95}{\micro\farad}
\end{itemize}

\subsubsection{Kondensator 3}

\paragraph{Marking code} Am Bauteil wurde der Wert \SI{10}{\micro\farad} abgelesen.

\paragraph{Programm} Vom Programm wurden folgende Werte abgelesen:

\begin{itemize}
    \item Charge Time: \SI{103276}{\micro\second}
    \item Kapazität: \SI{10.33}{\micro\farad}
\end{itemize}